% ---------------------------------------------------------------
% CHAPTER 7: CONCLUSIONS AND FUTURE WORKS
% ---------------------------------------------------------------
\chapter{Conclusions and Future Works}

\section{Conclusion}

This internship experience at IDFC FIRST Bank delivered an extensive, professional-grade 
computing engineering journey spanning all dimensions of contemporary cloud-based systems 
advancement. The structured pedagogical curriculum furnished comprehensive competencies in 
GoLang, examination-first programming, and quality source crafting. The SkyFox Cinemas 
iterative assignment furnished pragmatic experience in technology administration, 
interdisciplinary coordination, and incremental deployment under authentic time limitations.

The principal accomplishment encompasses the autonomous development and establishment of the 
\textbf{Robust Event Consumption Framework} --- a manufacturing-prepared GoLang utility that 
directly remedied a regulatory obligation in the bank's message stream-oriented service 
infrastructure. The framework's constituent subsystems --- reprocessing with geometric pause 
intervals and randomization, storage-mediated duplicate elimination, messaging zone distribution 
with incident management, utility-based information recovery, indicator compilation 
compatibility, and adaptable settings parameter --- combined accomplish complete data retention 
during brief system disruptions, 82\% evaluation protection, and library-style construction 
prepared for universal distribution via any information-oriented service at the bank.

The internship continues through June 2026. Remaining intervals are directed toward completion 
of in-progress SkyFox capabilities and targeted framework improvements realizable within the 
available calendar.

\section{Ongoing Work (May -- June 2026)}

\begin{figure}[H]
\centering
\begin{tikzpicture}[node distance=0.45cm]
  \node[rectangle, fill=idfcblue, text=white, minimum width=12.5cm,
        minimum height=0.8cm, font=\bfseries, rounded corners=4pt]
    (title) {Ongoing Work Plan --- May to June 2026};
  \node[widepurpleblock, below=0.4cm of title, text width=12.2cm,
        minimum height=1.1cm] (o1)
    {\textbf{Framework: Per-Topic Prometheus Metric Labels} ---
     Categorize currently present measurements and data series with Kafka transmitting 
     categorization and subscriber designation for individual-contributor measurement 
     without supplementary application programming};
  \node[widepurpleblock, below=0.4cm of o1, text width=12.2cm,
        minimum height=1.1cm] (o2)
    {\textbf{Framework: DLQ Replay Dry-Run Mode} ---
     Append a \texttt{--dry-run} parameter to the restorer utility facilitating technical 
     personnel to examine information prior to approving a mass-quantity restoration};
\end{tikzpicture}
\caption{Ongoing Work Plan (May -- June 2026)}
\end{figure}
\section{Future Works}

\subsection{Production Deployment and Environment Promotion}
The Robust Event Consumption Framework has received total construction, verification, and 
confirmation in a personal development context utilizing containerized transmission and data 
storage. The immediate subsequent phase is environmental advancement --- placement within 
the bank's examination context for synchronization verification towards genuine transmission 
structures, accompanied by supervised manufacturing initialization under technology 
operations direction. This encompasses creation of startup procedures, encrypted credential 
administration via Secrets management, and deployment robotics for programmatic assembly 
and readiness verification.

\subsection{Grafana SLO Dashboard}
The Prometheus measurements mandated for an operational metrics display --- ninety-ninth 
quantile information handling latency, message staging distribution percentage, reprocessing 
quantity distribution, and transmission handling lag --- are presently being maintained by 
the framework. The organic next activity encompasses constructing a monitoring screen that 
contrasts these opposing the bank's 99.9\% availability mandate, furnishing the critical 
procedures personnel a consolidated functional perspective of processing health spanning 
all implementations consuming the utilities.

\subsection{Kubernetes Deployment with Auto-Scaling}
The framework presently operates as a Docker deployment unit and underwent validation 
regionally employing Docker multi-container management. A scheduled improvement focuses on 
container orchestration positioning and customization of Pod Multiplication driven by 
transmission handling workload metrics via orchestration activation (Kubernetes Event-Driven 
Autoscaling). This capability would facilitate the consuming infrastructure to proportionally 
expand throughout peak transaction circumstances and proportionally contract throughout 
low-activity seasons, minimizing equipment provisioning without programmed management.

\subsection{OpenTelemetry Distributed Tracing}
The contemporary perceivability composition encompasses Prometheus measurement collection and 
formatted written records yet lacks distributed performance examination. Integrating the 
framework with telemetry abstraction (W3C Distribution Context specifications) would furnish 
end-to-end performance observation through the foundational transmission generator via the 
reprocessing mechanism to the ultimate following composition. This capability would accelerate 
performance disruption investigation across service perimeters substantially and constitutes a 
standard complement to the existent Prometheus coupling presently operational.

